\section*{الفصل \ref{ch:b2:basins-compared} --- مقارنة بين الأحواض الرسوبية}

\begin{solution}{exo:b2:basins-compared:1}
(أ) داخل كراتوني: هبوط بطيء على دفعات على مدى مئات ملايين السنين، دون تصدّع
كبير؛ وسببه موضع خلاف. (ب) شقّ: ترقيق ميكانيكي للقشرة أثناء التمطّط،
ثم انكماش حراري مع تبرّد الغلاف الصخري. (ج) هامش سلبي: الهبوط الحراري
الموروث من مرحلة الشقّ، يضخّمه حمل إسفين رسوبي سميك. (د) و(ه)
أمامي جبلي: الغلاف الصخري منثنٍ نحو الأسفل بحمل الحزام الدفعي، الألب والهيمالايا.
\end{solution}

\begin{solution}{exo:b2:basins-compared:2}
(أ) $Y = 3000\times(3330 - 2300)/(3330 - 1030) = 3000\times1030/2300 = \qty{1343}{m}$.
(ب) $1343 + 200 = \qty{1543}{m}$. (ج) دون حملها، كانت القاعدة ستقع على \qty{1343}{m}:
فالجزء الذي يفسح له الحمل المجال من السمك هو $1 - 1343/3000 = \qty{55}{\percent}$.
\end{solution}

\begin{solution}{exo:b2:basins-compared:3}
(أ) $\beta = 35/25 = 1.4$. (ب) $125/1.4 = \qty{89.3}{km}$. (ج) $S_i = 4334\times(1 - 1/1.4) =
\qty{1238}{m}$.
\end{solution}

\begin{solution}{exo:b2:basins-compared:4}
(أ) $\tau = (1.25\times10^{5})^2/(\pi^2\times8\times10^{-7}) = \qty{1.979e15}{s}$، أي
$1.979\times10^{15}/3.156\times10^{13} = \qty{62.7}{Myr}$. (ب) $62.7\times0.8 = \qty{50.2}{Myr}$.
(ج) $62.7\times(100/125)^2 = \qty{40.1}{Myr}$. (د) $1 - e^{-1} = \qty{63}{\percent}$؛ $1 -
e^{-3} = \qty{95}{\percent}$.
\end{solution}

\begin{solution}{exo:b2:basins-compared:5}
(أ) $0.63\,e^{-1000/1960} = 0.378$؛ $0.63\,e^{-2000/1960} = 0.227$. (ب) $T_s = 1000 -
0.63\times1960\times(0.6004 - 0.3605) = \qty{704}{m}$. (ج) $m = 2700\times704 + 1030\times296 =
\qty{2.206e6}{kg.m^{-2}}$؛ $\bar\rho = \qty{2206}{kg.m^{-3}}$.
\end{solution}

\begin{solution}{exo:b2:basins-compared:6}
(أ) $z_2' \leftarrow 704 + 1234.8\,(1 - e^{-z_2'/1960})$: 1077، 1226، 1278، 1296، 1301، 1303،
1304؛ فسمك الطبقة \qty{1304}{m}. (ب) $(2700\times704 + 1030\times600)/1304 =
\qty{1932}{kg.m^{-3}}$. (ج) بمقدار \qty{304}{m}، أي \qty{30}{\percent}: وهو الماء الذي
عصره الانضغاط.
\end{solution}

\begin{solution}{exo:b2:basins-compared:7}
(أ) $1 - e^{-100/62.7} = 0.7971$. من أجل $\beta = 1.5$: $3207\times0.4135\times0.7971 =
\qty{1057}{m}$؛ ومن أجل $\beta = 2$: $3207\times0.6366\times0.7971 = \qty{1627}{m}$. (ب) نعم: فالقيم من 1100 إلى
\qty{1400}{m} تقع بينهما. (ج) $\frac{\beta}{\pi}\sin\frac{\pi}{\beta} = 1250/(3207\times0.7971)
= 0.4890$، ومنه $\beta = 1.6 + 0.1\times(0.4890 - 0.4705)/(0.5204 - 0.4705) = 1.64$، ضمن
التمطّط من 50 إلى \qty{100}{\percent} المستنتج للحوض.
\end{solution}

\begin{solution}{exo:b2:basins-compared:8}
(أ) تحلّ كثافة الملء محلّ $\rho_w$ في المقام. (ب) البسط هو
$125\,000\times(1/3)\times74.69 = \num{3.112e6}$؛ $S_i = 3.112\times10^{6}/(3184.4 - 2300) =
\qty{3519}{m}$. (ج) تحت الماء، $3.112\times10^{6}/2154.4 = \qty{1445}{m}$: فالحوض المملوء بالرواسب
أعمق بمقدار 2.44 مرة.
\end{solution}

\begin{solution}{exo:b2:basins-compared:9}
(أ) $2171 - 100\times3330/2300 = 2171 - 145 = \qty{2026}{m}$. (ب) يتغيّر $Y$ مترًا بمتر
مع $W_d$: من 1971 إلى \qty{2371}{m}. (ج) عمق الماء، $\pm\qty{200}{m}$، مقابل
\qty{145}{m} لمستوى البحر.
\end{solution}

\begin{solution}{exo:b2:basins-compared:10}
(أ) إشارة $S_i$ هي إشارة القوس في بسطه، لأن المقام، $3184.4 - 1030$،
موجب؛ فيكون سالبًا، أي رفعًا، حين يكون القوس سالبًا. (ب) $(\rho_m - \rho_c)t_c/a =
\alpha T_1\rho_m/2$ يعطي $t_c = a\,\alpha T_1\rho_m/(2(\rho_m - \rho_c)) = 125\times0.04372\times
3330/1060 = \qty{17.17}{km}$. (ج) $t_c = 17.17/(1 - 0.04372\times17.17/250) = \qty{17.22}{km}$.
(د) حيث تكون القشرة أرقّ من نحو \qty{17}{km} قبل التمطّط: في الغلاف الصخري المحيطي،
أو في قشرة قارية رقّقها تصدّع سابق.
\end{solution}

\begin{solution}{exo:b2:basins-compared:11}
(أ) قبل التمطّط، يحوي العمود الغلاف الصخري ذا الكتلة $M = \rho_c t_c(1 - \alpha T_1
t_c/(2a)) + \rho_m(a - t_c)(1 - \alpha T_1(a + t_c)/(2a))$؛ وبعده، $\rho_w S_i + M/\beta +
\rho_m(1 - \alpha T_1)(a - S_i - a/\beta)$. وبمساواتهما، $S_i[\rho_m(1 - \alpha T_1) - \rho_w] =
(1 - 1/\beta)[\rho_m(1 - \alpha T_1)a - M]$، وبنشر $M$، $\rho_m(1 - \alpha T_1)a - M =
(\rho_m - \rho_c)t_c(1 - \alpha T_1 t_c/(2a)) - \rho_m\alpha T_1 a/2$: وهي الصيغة. (ب) المظهر النهائي هو $T_1 z/a$؛ ويزيد عليه المظهر الابتدائي
بمقدار $T_1 z(\beta - 1)/a$ فوق $a/\beta$ وبمقدار $T_1(1 - z/a)$ تحته. والتكاملان هما $T_1
a(\beta - 1)/(2\beta^2)$ و$T_1 a(\beta - 1)^2/(2\beta^2)$، ومجموعهما $T_1 a(1 - 1/\beta)/2$.
والانكماش، أي $\alpha$ مضروبًا في هذا، محمّلًا بالماء، يعطي العبارة. (ج) $3.28\times10^{-5}
\times3330\times1333\times125\,000\times0.5/(2\times2300) = \qty{1978}{m}$، مقابل $3207\times
2/\pi = \qty{2042}{m}$: فالحدّ الأول يبالغ في الهبوط النهائي بمقدار \qty{64}{m}، أي
\qty{3}{\percent}.
\end{solution}

\begin{solution}{exo:b2:basins-compared:12}
(أ) $3207\times0.6366\times(1 - e^{-150/62.7}) = 2042\times0.9086 = \qty{1855}{m}$. (ب) $1855\times
2300/(3330 - 2400) = \qty{4588}{m}$. (ج) من 6.9 إلى \qty{13.8}{km}. (د) يُغفل النموذج عدة
أمور قد تعمّق الحوض: غلاف صخري وشاحي مُطّ أكثر من القشرة، واستجابة إقليمية
(انثنائية) لحمل الرواسب، ورفع عابر للهامش أثناء التصدّع
يتبعه هبوط، أو إمداد بالرواسب كبير إلى حدّ أن الحوض امتلأ فوق طاقته وحُمّل
فوق مستوى سطح البحر على الرصيف.
\end{solution}

\begin{solution}{pb:b2:basins-compared:1}
\textbf{1.} $0.49\,e^{-4420/3704} = 0.149$ عند القمة؛ و$0.49\,e^{-6030/3704} = 0.096$ عند القاعدة.
\textbf{2.} $T_s = 1610 - 0.49\times3704\times(0.3032 - 0.1963) = \qty{1416}{m}$.
\textbf{3.} الطين الصفحي: $1110 - 0.63\times1960\times(e^{-3310/1960} - e^{-4420/1960}) = 1110 - 1234.8
\times(0.1847 - 0.1049) = \qty{1011}{m}$.
\textbf{4.} الرمل الطيني: $2680\times1004 + 1030\times696 = \qty{3.408e6}{kg.m^{-2}}$؛ والطباشير:
$2710\times1409 + 1030\times201 = \qty{4.025e6}{kg.m^{-2}}$؛ والطين الصفحي: $2700\times1011 + 1030\times99 =
\qty{2.832e6}{kg.m^{-2}}$؛ والرمل: $2650\times1416 + 1030\times194 = \qty{3.952e6}{kg.m^{-2}}$. والمجموع
\qty{1.4217e7}{kg.m^{-2}}، و$\bar\rho_s = 14\,216\,910/6030 = \qty{2358}{kg.m^{-3}}$.
\textbf{5.} $z_2' \leftarrow 1416 + 1815\,(1 - e^{-z_2'/3704})$: 1993، 2171، 2221، 2235، 2238، 2239،
2239: أي \qty{2239}{m}.
\textbf{6.} $(2650\times1416 + 1030\times823)/2239 = \qty{2055}{kg.m^{-3}}$.
\textbf{7.} $Y = 2239\times(3330 - 2055)/2300 + 200 = 1241 + 200 = \qty{1441}{m}$.
\textbf{8.} $Y = 6030\times(3330 - 2358)/2300 + 100 = 2548 + 100 = \qty{2648}{m}$. فالوحدات
مقيسة عند أعماقها الحالية: إنها منضغطة أصلًا بقدر انضغاطها اليوم.
\textbf{9.} $(2171 - 1441)/50 = \qty{14.6}{m}$ لكل مليون سنة من 150 إلى \qty{100}{Ma}؛
و$(2453 - 2171)/40 = \qty{7.05}{m}$ من 100 إلى \qty{60}{Ma}؛ و$(2648 - 2453)/60 = \qty{3.25}{m}$
منذ ذلك الحين: فالهبوط يتباطأ.
\textbf{10.} $6030 + 100 = \qty{6130}{m}$، أي 2.31 ضعف الهبوط التكتوني.
\textbf{11.} $(6130 - 2648)/6130 = \qty{57}{\percent}$.
\textbf{12.} $2171 - 100\times3330/2300 = \qty{2026}{m}$.
\textbf{13.} النقطة عند نهاية التصدّع، \qty{150}{Ma}: $S_i = \qty{1441}{m}$، إذا لم يحدث هبوط
حراري خلال \qty{15}{Myr} من التصدّع.
\textbf{14.} $\beta = 1/(1 - 1441/4334) = 1/0.6675 = 1.50$.
\textbf{15.} $35/1.50 = \qty{23.3}{km}$ من القشرة و$125/1.50 = \qty{83.3}{km}$ من الغلاف الصخري.
\textbf{16.} $(1.50/\pi)\sin(\pi/1.50) = 0.4775\times0.8660 = 0.4135$؛ $3207\times0.4135 =
\qty{1326}{m}$.
\textbf{17.} $S_i = 4334/3 = \qty{1445}{m}$. عند \qty{100}{Ma}، $t = \qty{50}{Myr}$: $1326\times(1 -
e^{-50/62.7}) = \qty{729}{m}$، $Y = \qty{2173}{m}$؛ وعند \qty{60}{Ma}: $1326\times0.7620 =
\qty{1010}{m}$، $Y = \qty{2455}{m}$؛ واليوم: $1326\times0.9086 = \qty{1205}{m}$، $Y = \qty{2650}{m}$.
\textbf{18.} النموذج أعمق بمقدار \qty{2}{m} عند كل نقطة: 2173 مقابل 2171، و2455 مقابل 2453،
و2650 مقابل 2648.
\textbf{19.} أعماق ماء قديمة لا تُعرف إلا في حدود مئة متر أو أكثر؛ ومستويات بحر في الماضي
مختلفة عن مستواه اليوم؛ وتصدّع دام \qty{15}{Myr}، فُقد خلاله جزء من الحرارة
أصلًا؛ والانثناء بدلًا من التوازن الأيزوستاتي المحلي؛ ووشاح مُطّ أكثر أو أقل من القشرة؛
وقشرة لم يكن سمكها \qty{35}{km} قبل التصدّع.
\textbf{20.} $1326\times(e^{-150/62.7} - e^{-200/62.7}) = \qty{67}{m}$.
\textbf{21.} $66.6\times2300/(3330 - 2400) = \qty{165}{m}$.
\textbf{22.} $1 - e^{-150/62.7} = \qty{91}{\percent}$.
\textbf{23.} $\tau\ln2 = \qty{43.5}{Myr}$ بعد نهاية التصدّع: أي عند \qty{106.5}{Ma}.
\textbf{24.} $\beta = \textbf{1.50}$: فقد مُطّت القشرة إلى مرة ونصف من عرضها.
\end{solution}
